% Luc Destombe --- licence CC BY-NC-SA 4.0
% https://creativecommons.org/licenses/by-nc-sa/4.0/deed.fr
% Fichier autonome : compiler avec pdflatex, deux fois (signets, nombre de pages).
\def\lyceeversion{1.13}\def\lyceerelectures{10}\def\lyceerodages{28}\def\lyceesceauannees{2026-2027}
\documentclass[11pt,a4paper]{article}
\makeatletter
\newif\ifeleve
\elevefalse

\newif\iflycee@palatino
\lycee@palatinotrue

\RequirePackage[utf8]{inputenc}
\RequirePackage[T1]{fontenc}
\RequirePackage[french]{babel}
\RequirePackage{etoolbox}

\newcommand{\ShorthandsOff}{\@ifundefined{shorthandoff}{}{\shorthandoff{;:!?}}}
\newcommand{\ShorthandsOn}{\@ifundefined{shorthandon}{}{\shorthandon{;:!?}}}
\ShorthandsOff
\AtBeginDocument{\ShorthandsOn}
\AtBeginEnvironment{tikzpicture}{\ShorthandsOff}

\RequirePackage{geometry}
\RequirePackage{fancyhdr}
\RequirePackage{lastpage}
\RequirePackage{multicol}
\RequirePackage{array}
\RequirePackage{enumitem}
\RequirePackage{xcolor}
\RequirePackage{needspace}

\RequirePackage{amsmath,amssymb}
\RequirePackage{mathpazo}
\renewcommand{\ttdefault}{lmtt}
\DeclareMathSizes{10.95}{10.95}{8}{5.5}
\begingroup\catcode`\_=\active \catcode`\^=\active
\gdef_#1{\sb{\scriptscriptstyle #1}}
\gdef^#1{\sp{\scriptscriptstyle\lycee@moinsepais #1}}
\endgroup
\newcommand\lycee@moins{\mathbin{%
  \setbox\z@\hbox{$\m@th\scriptscriptstyle\lycee@moinsorig$}%
  \dimen@=.55\wd\z@ \advance\dimen@-.5pt
  \hbox to.75\wd\z@{\hss$\m@th\scriptscriptstyle\vcenter{\hbox{%
    \tikz\draw[line width=.5pt,line cap=round](0,0)--(\the\dimen@,0);}}$\hss}}}
\begingroup\lccode`\~=`\- \lowercase{\endgroup\let~}\lycee@moins
\newcommand\lycee@moinsepais{\mathcode`\-="8000 }
\AtBeginDocument{\mathchardef\lycee@moinsorig=\mathcode`\-\relax}
\AtBeginDocument{\catcode`\_=12 \mathcode`\_="8000
  \catcode`\^=12 \mathcode`\^="8000 }
\newcommand\lycee@exposants{%
  \fontdimen13\textfont2=.48\fontdimen6\textfont2
  \fontdimen14\textfont2=.48\fontdimen6\textfont2
  \fontdimen15\textfont2=.48\fontdimen6\textfont2
  \fontdimen13\scriptfont2=.48\fontdimen6\scriptfont2
  \fontdimen14\scriptfont2=.48\fontdimen6\scriptfont2
  \fontdimen15\scriptfont2=.48\fontdimen6\scriptfont2}
\every@math@size\expandafter{\the\every@math@size\lycee@exposants}
\newlength{\retraitcalculs}
\setlength{\retraitcalculs}{2em}

\RequirePackage{tikz}
\usetikzlibrary{arrows.meta,calc,babel,shapes.misc}
\RequirePackage[most]{tcolorbox}
\tcbuselibrary{skins,breakable}

\definecolor{coulDef}{HTML}{1F4E79}
\definecolor{coulProp}{HTML}{7030A0}
\definecolor{coulRem}{HTML}{C55A11}
\definecolor{coulEx}{HTML}{2E7D32}
\definecolor{coulExo}{HTML}{B03A2E}
\definecolor{coulPart}{HTML}{1F4E79}
\definecolor{coulMeth}{HTML}{1B6E4F}
\definecolor{coulCourbe}{HTML}{0B5ED7}
\definecolor{coulCourbeB}{HTML}{C00000}
\definecolor{coulCourbeC}{HTML}{2E7D32}
\definecolor{coulGrille}{HTML}{8C8C8C}
\definecolor{coulRep}{HTML}{1B6E4F}
\definecolor{coulBar}{HTML}{2E7D32}

\newcommand{\R}{\mathbb{R}}
\newcommand{\Rs}{\mathbb{R}^{*}}
\renewcommand{\le}{\leqslant}
\renewcommand{\ge}{\geqslant}
\renewcommand{\approx}{\simeq}
\newcommand{\vect}[1]{\overrightarrow{#1}}
\newcommand{\norme}[1]{\left\lVert #1 \right\rVert}
\newcommand{\intervalle}[2]{\left[#1\,;#2\right]}
\RequirePackage{cancel}
\renewcommand{\CancelColor}{\color{coulExo}}
\newcommand{\fois}[1]{\times{\color{coulExo}#1}}
\newcommand{\barre}[1]{\cancel{#1}}

\tikzset{grille/.style={coulGrille,line width=0.4pt}}
\newcommand{\quadrillage}[4]{\draw[grille,step=1] (#1,#2) grid (#3,#4);}
\tikzset{
  croix/.style={cross out,draw,line width=0.6pt,inner sep=0pt,minimum size=4pt},
  croixdroite/.style={croix,rotate=45,line width=0.8pt},
}
\newcommand{\pt}[3]{\path (#1) node[croix]{} node[#2,font=\small,inner sep=2.5pt]{$#3$};}

\newcommand{\lycee@repere}[4]{%
  \draw[grille,step=1] (#1,#2) grid (#3,#4);
  \draw[-{Stealth[length=2mm]},thick] (#1,0)--({#3+0.4},0) node[below left=-1pt]{$x$};
  \draw[-{Stealth[length=2mm]},thick] (0,#2)--(0,{#4+0.4}) node[below left=-1pt]{$y$};
  \node[below left,font=\scriptsize,inner sep=1.5pt] at (0,0) {$0$};}
\newcommand{\lycee@gradx}[1]{\foreach \k/\l in {#1}{%
  \draw (\k,0.07)--(\k,-0.07) node[below,font=\scriptsize,inner sep=1.5pt]{$\l$};}}
\newcommand{\lycee@grady}[1]{\foreach \k/\l in {#1}{%
  \draw (0.07,\k)--(-0.07,\k) node[left,font=\scriptsize,inner sep=1.5pt]{$\l$};}}
\newcommand{\lycee@intff}[2]{\left[#1\,;#2\right]}
\newcommand{\lycee@intfo}[2]{\left[#1\,;#2\right[}
\newcommand{\lycee@intof}[2]{\left]#1\,;#2\right]}
\newcommand{\lycee@intoo}[2]{\left]#1\,;#2\right[}
\newcommand{\lycee@ens}[1]{\left\{#1\right\}}
\AtBeginDocument{%
  \providecommand{\repere}{\lycee@repere}%
  \providecommand{\gradx}{\lycee@gradx}%
  \providecommand{\grady}{\lycee@grady}%
  \providecommand{\Cf}{\mathcal{C}\sb{\scriptscriptstyle f}}%
  \providecommand{\Cg}{\mathcal{C}\sb{\scriptscriptstyle g}}%
  \providecommand{\intff}{\lycee@intff}%
  \providecommand{\intfo}{\lycee@intfo}%
  \providecommand{\intof}{\lycee@intof}%
  \providecommand{\intoo}{\lycee@intoo}%
  \providecommand{\ens}{\lycee@ens}}

\newlist{questions}{enumerate}{3}
\setlist[questions,1]{label=\arabic*\textdegree),leftmargin=*,itemsep=2pt,topsep=3pt}
\setlist[questions,2]{label=\alph*),leftmargin=*,itemsep=1pt,topsep=2pt}
\setlist[questions,3]{label=\roman*),leftmargin=*,itemsep=1pt,topsep=1pt}
\setlist[itemize]{label=\textbullet,leftmargin=*,itemsep=2pt,topsep=3pt}

\newcommand{\besoinplace}[1]{\par\penalty\@M\begingroup
  \dimen@=#1\relax\dimen@ii=\pagegoal\advance\dimen@ii-\pagetotal
  \ifdim\dimen@>\dimen@ii\ifdim\dimen@ii>\z@\vfil\fi\break\fi\endgroup}

\newcommand{\lycee@pied}{\thepage/\pageref{LastPage}}
\newcommand{\entete}[2]{%
  \pagestyle{fancy}\fancyhf{}%
  \fancyhead[L]{\footnotesize\color{coulPart}\textbf{#1}}%
  \fancyhead[R]{\footnotesize\color{coulPart}\textbf{#2}}%
  \fancyfoot[C]{\footnotesize\lycee@pied}%
  \renewcommand{\headrulewidth}{0.4pt}%
  \renewcommand{\headrule}{\hbox to\headwidth{%
    \color{coulPart!40}\leaders\hrule height \headrulewidth\hfill}}}

\RequirePackage{ccicons}
\newcommand{\lycee@licence}{{\scriptsize\color{gray}%
  \copyright\ Luc Destombe --- \ccbyncsa\ CC BY-NC-SA 4.0}}
\newif\iflycee@licence \lycee@licencetrue
\newcommand{\sanslicence}{\lycee@licencefalse}
\AtBeginDocument{\iflycee@licence\AddToHookNext{shipout/foreground}{%
  \put(0,0){\rlap{\kern\dimexpr1in+\hoffset+\oddsidemargin+\textwidth\relax
    \llap{\raisebox{-\dimexpr1in+\voffset+\topmargin+\headheight+\headsep
      +\textheight+\footskip\relax}{\lycee@licence}}}}}\fi}

\newcommand{\lycee@site}{faire-des-maths.fr}
\newcommand{\lycee@version}{{\scriptsize\color{gray}\lycee@site
  \ifdefined\lyceeversion\ --- Version~\lyceeversion\fi}}
\AtBeginDocument{\iflycee@licence\AddToHookNext{shipout/foreground}{%
  \put(0,0){\rlap{\kern\dimexpr1in+\hoffset+\oddsidemargin\relax
    \raisebox{-\dimexpr1in+\voffset+\topmargin+\headheight+\headsep
      +\textheight+\footskip\relax}{\lycee@version}}}}\fi}

\newcommand{\formulesserrees}{%
  \setlength{\abovedisplayskip}{3pt}\setlength{\belowdisplayskip}{3pt}%
  \setlength{\abovedisplayshortskip}{2pt}\setlength{\belowdisplayshortskip}{3pt}}

\setlength{\parindent}{0pt}

\RequirePackage[implicit=false,hidelinks,pdfencoding=auto,psdextra,bookmarksopen,
  bookmarksopenlevel=1,pdfcreator={},pdfproducer={}]{hyperref}
\RequirePackage{bookmark}
\hypersetup{pdfauthor={Luc Destombe},pdfkeywords={CC BY-NC-SA 4.0}}
\newcounter{lycee@signet}
\newcommand{\signet}[3][]{\stepcounter{lycee@signet}%
  \edef\lycee@dest{\if\relax\detokenize{#1}\relax
    signet.\arabic{lycee@signet}\else#1\fi}%
  \hypertarget{\lycee@dest}{}\bookmark[level=#2,dest=\lycee@dest]{#3}}
\pdfstringdefDisableCommands{%
  \def\R{R}\def\vect#1{#1}\def\Cf{Cf}\def\Cg{Cg}\def\fois#1{×#1}%
  \def\barre#1{#1}\def\intervalle#1#2{[#1;#2]}\def\intff#1#2{[#1;#2]}%
  \def\textsuperscript#1{#1}\def\dfrac#1#2{#1/#2}\def\frac#1#2{#1/#2}%
  \def\vec#1{#1⃗}}%
\newcommand{\res}[1]{#1}
\geometry{top=2cm,bottom=1cm,left=1cm,right=1cm,headheight=14pt,footskip=0.6cm}

\RequirePackage{pgfplots}
\pgfplotsset{compat=1.18}
\RequirePackage{tkz-tab}

\newcounter{partiecours}
\renewcommand{\thepartiecours}{\Roman{partiecours}}
\newcommand{\partiecours}[1]{%
  \refstepcounter{partiecours}%
  \Needspace*{8\baselineskip}%
  \bigskip
  \noindent\begin{tcolorbox}[enhanced,colback=coulPart,colframe=coulPart,
      boxrule=0pt,arc=2pt,left=8pt,right=8pt,top=4pt,bottom=4pt,
      before skip=6pt,after skip=10pt]
    \color{white}\bfseries\large \thepartiecours{} --- #1%
    \signet[partie-\arabic{partiecours}]{1}{\thepartiecours{} --- #1}
  \end{tcolorbox}%
  \addcontentsline{toc}{section}{\thepartiecours{} --- #1}%
}

\newtcolorbox{boitecours}[3][]{%
  enhanced,breakable,
  colback=white,colframe=#2,coltitle=white,
  fonttitle=\bfseries,
  attach boxed title to top left={xshift=8pt,yshift=-\tcboxedtitleheight/2},
  boxed title style={colback=#2,arc=2pt,boxrule=0pt},
  boxrule=0.9pt,arc=2pt,
  left=8pt,right=8pt,top=8pt,bottom=6pt,
  title={#3},#1}

\newcounter{methode}
\newenvironment{definitionbox}[1][Définition]%
  {\begin{boitecours}[phantom={\signet{2}{#1}}]{coulDef}{#1}}{\end{boitecours}}
\newenvironment{proprietebox}[1][Propriété]%
  {\begin{boitecours}[phantom={\signet{2}{#1}}]{coulProp}{#1}}{\end{boitecours}}
\newenvironment{methodebox}[1][Méthode]{\stepcounter{methode}%
  \begin{boitecours}[phantom={\signet[methode-\arabic{methode}]{2}{#1}}]{coulMeth}{#1}}%
  {\end{boitecours}}

\newcommand{\etape}[1]{\par\smallskip\textbf{\color{coulMeth}#1}\par\nobreak\smallskip}
\newcommand{\enonce}[1]{\emph{#1}\par\smallskip}
\newcommand{\reponse}[1]{\par\smallskip
  {\footnotesize\itshape\color{black!55}Réponses : #1}}
\newcommand{\demo}[1]{\textbf{\color{coulMeth}Démonstration.} #1}
\newcommand{\afaire}[1]{%
  \par\smallskip
  \noindent{\small\color{coulExo}$\blacktriangleright$\ \textbf{À faire :}
    fiche d'exercices du chapitre, #1.}\par\smallskip}

\newenvironment{calculs}{\setlength{\jot}{5pt}\@fleqntrue\@mathmargin\retraitcalculs
  \setlength{\abovedisplayskip}{4pt}\setlength{\belowdisplayskip}{4pt}%
  \setlength{\abovedisplayshortskip}{4pt}\setlength{\belowdisplayshortskip}{4pt}%
  \csname alignat*\endcsname{2}}%
  {\csname endalignat*\endcsname}
\newcommand{\rem}[1]{\text{\small\itshape\color{black!60}#1}}
\newcommand{\explique}[2]{\par\smallskip\noindent
  \begin{minipage}[t]{0.57\linewidth}\raggedright #1\end{minipage}\hfill
  \begin{minipage}[t]{0.40\linewidth}\raggedright\leavevmode
    {\small\itshape\color{black!60}#2}\end{minipage}\par}
\newcommand{\cas}[1]{\par\smallskip\noindent\textbf{\color{coulExo}#1}\vadjust{\nobreak}}
\newcommand{\calc}[1]{\par\smallskip\textbf{\color{coulExo}#1}\par\nobreak}

\newtcolorbox{remarquebox}[1][Remarque]{%
  enhanced,breakable,colback=white,colframe=coulRem,
  boxrule=0pt,leftrule=2.5pt,arc=0pt,outer arc=0pt,
  left=8pt,right=6pt,top=5pt,bottom=5pt,
  before upper={\textbf{\color{coulRem}#1}\par\smallskip}}

\newtcolorbox{exemplebox}[1][Exemple]{%
  enhanced,breakable,colback=white,colframe=coulEx,
  boxrule=0pt,leftrule=2.5pt,arc=0pt,outer arc=0pt,
  left=8pt,right=6pt,top=5pt,bottom=5pt,
  before upper={\textbf{\color{coulEx}#1}\par\smallskip}}

\pgfplotsset{
  repere/.style={
    axis lines=middle,
    axis line style={-{Stealth[length=2.4mm]},thick},
    every axis x label/.style={at={(ticklabel* cs:1.0)},anchor=west},
    every axis y label/.style={at={(ticklabel* cs:1.0)},anchor=south},
    label style={font=\footnotesize},
    tick label style={font=\scriptsize},
    grid=both,
    major grid style={grille},
    minor tick num=0,
    tick align=inside,
    clip mode=individual,
  },
}
\tikzset{
  courbe/.style={coulCourbe,line width=1.1pt,smooth},
  courbeB/.style={coulCourbeB,line width=1.1pt,smooth},
  courbeC/.style={coulCourbeC,line width=1.1pt,smooth},
}

\newlist{etapes}{enumerate}{1}
\setlist[etapes]{label=\textbf{\arabic*.},leftmargin=*,itemsep=1pt,topsep=2pt}

\newlist{expressions}{enumerate}{1}
\setlist[expressions]{label={\Alph*\ $=$},leftmargin=*,itemsep=3pt,topsep=3pt}

\newcounter{exo}
\renewcommand{\theexo}{\ifnum\value{exo}<10 0\fi\arabic{exo}}
\newtcolorbox{exercice}[1][]{%
  enhanced,breakable,
  colback=white,colframe=coulExo,
  boxrule=0.9pt,arc=2pt,
  left=8pt,right=8pt,top=8pt,bottom=6pt,
  coltitle=white,fonttitle=\bfseries,
  attach boxed title to top left={xshift=8pt,yshift=-\tcboxedtitleheight/2},
  boxed title style={colback=coulExo,arc=2pt,boxrule=0pt},
  phantom={\signet{2}{Exercice \arabic{exo}}},
  title={Exercice \theexo\ifx\relax#1\relax\else{} --- #1\fi}}
\newenvironment{exo}[1][\relax]{\refstepcounter{exo}\begin{exercice}[#1]}{\end{exercice}}

  \newenvironment{correction}{\par}{\par}
  \newcommand{\autableau}{}
  \newcommand{\etapeexemples}{\etape{Exemples corrigés}}
  \newcommand{\etapetableau}[1]{\etape{#1}}
  \newenvironment{exempletableau}[1][Exemple]%
    {\begin{exemplebox}[#1]}{\end{exemplebox}}

\RequirePackage{comment}
\excludecomment{correctionavj}

\newcommand{\coursEntete}{}
\newcommand{\coursNiveau}{}
\newcommand{\coursTitre}{}
\newcommand{\coursSurTitre}{}
\newcommand{\coursSousTitre}{}

\newcommand{\enteteCours}[1]{\renewcommand{\coursEntete}{#1}}
\newcommand{\chapitrecours}[2]{\enteteCours{Chapitre #1 --- #2}}
\newcommand{\niveaucours}[1]{\renewcommand{\coursNiveau}{#1}}
\newcommand{\titrecours}[3][]{%
  \renewcommand{\coursTitre}{#2}%
  \renewcommand{\coursSurTitre}{#1}%
  \renewcommand{\coursSousTitre}{#3}}

\pagestyle{fancy}
\fancyhf{}
\fancyhead[L]{\footnotesize\color{coulPart}\textbf{\coursEntete}}
\fancyhead[R]{\footnotesize\color{coulPart}\textbf{\coursNiveau}}
\fancyfoot[C]{\footnotesize\thepage}
\renewcommand{\headrulewidth}{0.4pt}
\renewcommand{\headrule}{\hbox to\headwidth{\color{coulPart!40}\leaders\hrule height \headrulewidth\hfill}}

\setlength{\parskip}{2pt}

\AtBeginDocument{%
  \begin{center}
    \begin{tcolorbox}[enhanced,width=0.72\linewidth,colback=white,colframe=coulPart,
        boxrule=1.6pt,arc=3pt,halign=center,top=10pt,bottom=10pt]
      \ifx\coursSurTitre\empty
        {\Huge\bfseries\color{coulPart} \coursTitre}\\[4pt]
      \else
        {\Huge\bfseries\color{coulPart} \coursTitre}\\[3pt]
        {\Large\bfseries\color{coulPart} \coursSurTitre}\\[5pt]
      \fi
      {\normalsize \coursSousTitre}%
      
    \end{tcolorbox}
  \end{center}%
}

\usetikzlibrary{decorations.text}
\definecolor{lycee@metal}{HTML}{1F4E79}
\definecolor{lycee@manche}{HTML}{8A5A2B}
\definecolor{lycee@etincelle}{HTML}{F28C28}
\newcommand{\lycee@pluriel}[2]{#1 #2\ifnum#1>1 s\fi}
\newcommand{\lycee@sceau}{%
  \edef\lycee@bas{\lycee@pluriel{\lyceerelectures}{relecture} et
    \lycee@pluriel{\lyceerodages}{rodage}}%
  \begin{tikzpicture}[line join=round]
    \fill[coulPart!10,draw=coulPart,line width=.8pt] (0:2.2)
      \foreach \i in {1,...,71}{ -- (\i*5:{isodd(\i)?2.1:2.2}) } -- cycle;
    \draw[coulPart,line width=1.2pt] (0,0) circle (1.93);
    \draw[coulPart,line width=.4pt] (0,0) circle (1.85);
    \draw[coulPart,line width=.4pt] (0,0) circle (1.2);
    \path[decorate,decoration={text along path,text align=center,
        text={|\fontsize{9.5}{11}\selectfont\bfseries\color{coulPart}|TEST{\'E} AVEC MES {\'E}L{\`E}VES}}]
      (180:1.45) arc (180:0:1.45);
    \path[decorate,decoration={text along path,text align=center,
        text/.expanded={|\noexpand\fontsize{8.5}{10}\noexpand\selectfont\noexpand\color{coulPart}|\lycee@bas}}]
      (195:1.68) arc (195:345:1.68);
    \begin{scope}[x=.25mm,y=-.25mm]
      \begin{scope}[shift={(-32,-40)}]%
        \fill[lycee@metal] (23,31) rectangle (41,49);
        \draw[lycee@etincelle,line width=.5pt,line cap=round]
          (20,30)--(15,27) (19,34)--(14,35) (44,30)--(49,27) (45,34)--(50,35);
        \fill[lycee@manche,rounded corners=.3mm] (38,24.6) rectangle (60,27.6);
        \fill[lycee@metal,rounded corners=.4mm] (25,20) rectangle (39,31);
        \fill[lycee@etincelle,rounded corners=.2mm] (25,20) rectangle (28,31);
      \end{scope}
    \end{scope}
    \ifdefined\lyceesceauplusieurs\def\lycee@annee{Ann{\'e}es~:}\else
      \def\lycee@annee{Ann{\'e}e}\fi
    \path[decorate,decoration={text along path,text align=center,
        text/.expanded={|\noexpand\fontsize{8}{10}\noexpand\selectfont\noexpand\color{coulPart}|%
          \unexpanded\expandafter{\lycee@annee} \unexpanded\expandafter{\lyceesceauannees}}}]
      (95:1.08) arc (95:445:1.08);
  \end{tikzpicture}}
\AtEndDocument{\ifdefined\lyceesceauannees
  \par\besoinplace{5cm}\vspace{1.5em}%
  \begin{center}\lycee@sceau\end{center}\fi}
\makeatother

\chapitrecours{3}{Vecteurs}
\niveaucours{Seconde}
\titrecours{VECTEURS}{Cours --- Seconde}

\begin{document}

\ShorthandsOff
\tikzset{
  vecteur/.style={-{Stealth[length=2.6mm]},line width=1.1pt,coulCourbe},
  vecteurB/.style={-{Stealth[length=2.6mm]},line width=1.1pt,coulCourbeB},
  vecteurC/.style={-{Stealth[length=2.6mm]},line width=1.1pt,coulCourbeC},
  aide/.style={dashed,black!45,line width=0.6pt},
  fig/.style={line width=0.8pt,black!75},
}
\ShorthandsOn

\begin{remarquebox}[Comment utiliser ce cours]
Chaque partie commence par ce qu'il faut \textbf{savoir} (définitions, propriétés), puis
vient ce qu'il faut \textbf{savoir faire} : une méthode, avec des
\textbf{exemples corrigés} faits en classe, puis des exercices
\og\textbf{À vous de jouer}\fg{} à chercher seul. Les figures se font au crayon et à la
règle : un vecteur mal tracé est un vecteur faux.

\end{remarquebox}

\partiecours{Caractéristiques d'un vecteur}

\begin{definitionbox}[Définition 1 --- Vecteur]
\begin{minipage}[t]{0.58\linewidth}\raggedright
Deux points $A$ et $B$ définissent le \textbf{vecteur} $\vect{AB}$ : il va de son
\textbf{origine} $A$ à son \textbf{extrémité} $B$. On le représente par une flèche
de $A$ vers $B$. Il a trois caractéristiques :
\begin{itemize}
  \item une \textbf{direction}, celle de la droite $(AB)$ ;
  \item un \textbf{sens}, de $A$ vers $B$ ;
  \item une \textbf{longueur}, appelée \textbf{norme} et notée $\norme{\vect{AB}}=AB$.
\end{itemize}
On nomme aussi un vecteur par une seule lettre : $\vec{u}$, $\vec{v}$\dots
\end{minipage}\hfill
\begin{minipage}[t]{0.38\linewidth}
\vspace{0pt}\centering
\begin{tikzpicture}[scale=0.8]
  \draw[aide] (-1,-0.5) -- (5,2.5) node[right,font=\footnotesize,black]{$(AB)$};
  \draw[vecteur] (0,0) -- node[midway,above left,font=\footnotesize]{$\vect{AB}$} (4,2);
  \pt{0,0}{below right}{A} \pt{4,2}{below right}{B}
  \node[font=\footnotesize,below=0.35cm] at (0,0) {origine};
  \node[font=\footnotesize,below=0.35cm] at (4,2) {extrémité};
\end{tikzpicture}
\end{minipage}
\end{definitionbox}

\begin{remarquebox}[Deux confusions à éviter]
\begin{itemize}
  \item La \textbf{direction} n'est pas le \textbf{sens} : une direction a deux sens.
        $\vect{AB}$ et $\vect{BA}$ ont la même direction et des sens contraires.
  \item $AB$ est un nombre (une distance), $\vect{AB}$ est un vecteur : on écrit
        $\norme{\vect{AB}}=3$, jamais $\vect{AB}=3$.
\end{itemize}
\end{remarquebox}

\begin{definitionbox}[Définition 2 --- Vecteurs égaux, nul, opposés]
\begin{itemize}
  \item Deux vecteurs sont \textbf{égaux} s'ils ont même direction, même sens et même
        norme. Un vecteur peut donc être copié et déplacé : on dit qu'il a une infinité
        de \textbf{représentants}.
  \item Le vecteur $\vect{AA}$ est le \textbf{vecteur nul}, noté $\vec{0}$.
  \item Deux vecteurs sont \textbf{opposés} s'ils ont même direction, même norme et
        des sens contraires : $\vect{BA}=-\vect{AB}$.
\end{itemize}
\end{definitionbox}

\begin{exemplebox}[Exemple --- Trois représentants d'un même vecteur]
\begin{center}
\begin{tikzpicture}[scale=0.55]
  \quadrillage{-0.4}{-0.4}{11.4}{3.4}
  \draw[vecteur] (0,0) -- (3,1);  \node[font=\small,coulCourbe] at (1.2,1) {$\vec{u}$};
  \draw[vecteur] (4,2) -- (7,3);  \node[font=\small,coulCourbe] at (5.2,3) {$\vec{u}$};
  \draw[vecteur] (8,0) -- (11,1); \node[font=\small,coulCourbe] at (9.2,1) {$\vec{u}$};
\end{tikzpicture}
\end{center}
Les trois flèches font le même déplacement ($3$ carreaux à droite et $1$ vers le haut) :
ce sont trois représentants du même vecteur $\vec{u}$.
\end{exemplebox}

\begin{methodebox}[Méthode 1 --- Reconnaître des vecteurs égaux ou opposés]
Sur un quadrillage, on lit le \textbf{déplacement} de chaque flèche : tant de carreaux
horizontalement, tant de carreaux verticalement.
\etapeexemples
\begin{questions}
  \item Tracer les vecteurs $\vect{AB}$, $\vect{CD}$, $\vect{EF}$ et $\vect{GH}$.
  \item Pour chaque vecteur $\vect{CD}$, $\vect{EF}$ et $\vect{GH}$, dire s'il est égal à
        $\vect{AB}$, opposé à $\vect{AB}$, ou ni l'un ni l'autre.
  \item Comparer $\vect{FE}$ et $\vect{AB}$.
\end{questions}
\begin{center}
\begin{tikzpicture}[scale=0.6]
  \quadrillage{-0.4}{-0.4}{10.4}{3.4}
  \pt{0,0}{below left}{A} \pt{4,1}{below right}{B}
  \pt{1,2}{above left}{C} \pt{5,3}{above right}{D}
  \pt{10,3}{above right}{E} \pt{6,2}{above left}{F}
  \pt{7,0}{below left}{G} \pt{10,1}{below right}{H}
\end{tikzpicture}
\end{center}
\begin{correction}
\cas{1\textdegree)}\par\nopagebreak
\begin{center}
\begin{tikzpicture}[scale=0.6]
  \quadrillage{-0.4}{-0.4}{10.4}{3.4}
  \draw[vecteur] (0,0) -- node[midway,above,font=\footnotesize]{$\vect{AB}$} (4,1);
  \draw[vecteur] (1,2) -- node[midway,above,font=\footnotesize]{$\vect{CD}$} (5,3);
  \draw[vecteurB] (10,3) -- node[midway,above,font=\footnotesize]{$\vect{EF}$} (6,2);
  \draw[vecteurC] (7,0) -- node[midway,above,font=\footnotesize]{$\vect{GH}$} (10,1);
  \pt{0,0}{below left}{A} \pt{4,1}{below right}{B}
  \pt{1,2}{above left}{C} \pt{5,3}{above right}{D}
  \pt{10,3}{above right}{E} \pt{6,2}{above left}{F}
  \pt{7,0}{below left}{G} \pt{10,1}{below right}{H}
\end{tikzpicture}
\end{center}
\explique{Pour $\vect{AB}$, on trace une flèche qui part de $A$ et dont la pointe est en
$B$.}{l'origine est la première lettre}
\explique{De même, $\vect{CD}$ va de $C$ à $D$, $\vect{EF}$ de $E$ à $F$, $\vect{GH}$ de
$G$ à $H$.}{la pointe est sur la seconde lettre}
\cas{2\textdegree)}\par\nopagebreak
\explique{$\vect{AB}$ : $4$ carreaux à droite et $1$ vers le haut.}{le vecteur de
référence}
\explique{$\vect{CD}$ : $4$ à droite, $1$ vers le haut, donc $\vect{CD}=\vect{AB}$.}{même
déplacement}
\explique{$\vect{EF}$ : $4$ à gauche, $1$ vers le bas, donc $\vect{EF}=-\vect{AB}$.}{déplacement
contraire}
\explique{$\vect{GH}$ : $3$ à droite, $1$ vers le haut : ni égal, ni opposé.}{autre
déplacement}
\cas{3\textdegree)}
\explique{$\vect{FE}$ : de $F$ à $E$, $4$ carreaux à droite et $1$ vers le haut, comme
$\vect{AB}$ : $\vect{FE}=\vect{AB}$.}{échanger les lettres de $\vect{EF}$ change le sens}
\end{correction}

\etape{À vous de jouer}
\begin{center}
\begin{tikzpicture}[scale=0.6]
  \quadrillage{-0.4}{-0.4}{11.4}{4.4}
  \pt{0,0}{below left}{M} \pt{3,2}{above left}{N}
  \pt{4,0}{below right}{V} \pt{7,2}{right}{W}
  \pt{7,4}{above right}{R} \pt{4,2}{above left}{S}
  \pt{8,0}{below left}{T} \pt{11,1}{below right}{U}
\end{tikzpicture}
\end{center}
\begin{questions}
  \item Tracer les vecteurs $\vect{MN}$, $\vect{VW}$, $\vect{RS}$ et $\vect{TU}$.
  \item Pour chaque vecteur $\vect{VW}$, $\vect{RS}$ et $\vect{TU}$, dire s'il est égal à
        $\vect{MN}$, opposé à $\vect{MN}$, ou ni l'un ni l'autre.
  \item Comparer $\vect{SR}$ et $\vect{MN}$.
\end{questions}
\reponse{2\textdegree) $\vect{VW}=\vect{MN}$, $\vect{RS}=-\vect{MN}$, $\vect{TU}$ ni l'un
ni l'autre.\; 3\textdegree) $\vect{SR}=\vect{MN}$ : échanger les lettres change le sens.}
\end{methodebox}

\afaire{exercices 1 à 4}

\partiecours{Vecteurs et déplacements dans un quadrillage}

\begin{definitionbox}[Définition 3 --- Déplacement d'un vecteur]
\begin{minipage}[t]{0.58\linewidth}\raggedright
Sur un quadrillage, un vecteur se décrit par son \textbf{déplacement}, compté en
carreaux de l'origine vers l'extrémité :
\begin{itemize}
  \item d'abord \textbf{horizontalement} : tant de carreaux à droite ou à gauche ;
  \item puis \textbf{verticalement} : tant de carreaux vers le haut ou vers le bas.
\end{itemize}
Deux vecteurs égaux ont le même déplacement.

\smallskip
Deux vecteurs opposés ont des déplacements \textbf{contraires} : ci-contre, $\vec{u}$ fait
$4$ carreaux à droite et $3$ vers le haut ; son opposé $-\vec{u}$ fait $4$ carreaux à
gauche et $3$ vers le bas.
\end{minipage}\hfill
\begin{minipage}[t]{0.38\linewidth}
\vspace{0pt}\centering
\begin{tikzpicture}[scale=0.55]
  \quadrillage{-0.4}{-6.4}{5.4}{3.4}
  \draw[aide,black!70] (0,0) -- node[midway,below,font=\footnotesize,black]{$4$ à droite} (4,0);
  \draw[aide,black!70] (4,0) -- node[midway,right,font=\footnotesize,black]{$3$ vers le haut} (4,3);
  \draw[vecteur] (0,0) -- node[midway,above left,font=\footnotesize]{$\vec{u}$} (4,3);
  \draw[aide,black!70] (4,-3) -- node[midway,above,font=\footnotesize,black]{$4$ à gauche} (0,-3);
  \draw[aide,black!70] (0,-3) -- node[midway,left,align=right,font=\footnotesize,black]{$3$ vers\\le bas} (0,-6);
  \draw[vecteurB] (4,-3) -- node[midway,below right,font=\footnotesize]{$-\vec{u}$} (0,-6);
\end{tikzpicture}
\end{minipage}
\end{definitionbox}

\begin{methodebox}[Méthode 2 --- Lire et tracer un déplacement]
Pour \textbf{lire}, on compte les carreaux de l'origine à la pointe. Pour \textbf{tracer}
depuis une origine, on fait le déplacement et on place la pointe. Pour tracer un vecteur
qui doit \textbf{arriver} en un point, on fait le déplacement à l'envers depuis ce point :
on trouve l'origine. Pour tracer l'\textbf{opposé} $-\vec{u}$, on fait le déplacement
contraire de celui de $\vec{u}$.
\etapeexemples
\begin{questions}[label=\alph*)]
  \item Lire le déplacement de chaque vecteur $\vec{u}$, $\vec{v}$ et $\vec{w}$.
  \item Tracer le vecteur $\vec{s}$ d'origine $A$ qui fait $4$ carreaux à droite et $1$
        vers le bas.
  \item Tracer un représentant de $\vec{u}$ d'origine $B$.
  \item Tracer un autre représentant de $\vec{u}$, d'extrémité $C$.
  \item Tracer un représentant de $-\vec{u}$ d'origine $D$.
\end{questions}
\begin{center}
\begin{tikzpicture}[scale=0.45]
  \node[font=\small\bfseries,coulExo] at (7,-1.3) {a)};
  \quadrillage{-0.4}{-0.4}{14.4}{4.4}
  \draw[vecteur] (0,1) -- node[midway,above left,font=\footnotesize]{$\vec{u}$} (4,3);
  \draw[vecteurB] (6,4) -- node[midway,right,font=\footnotesize]{$\vec{v}$} (9,0);
  \draw[vecteurC] (14,3) -- node[midway,above left,font=\footnotesize]{$\vec{w}$} (11,1);
\end{tikzpicture}
\quad
\begin{tikzpicture}[scale=0.45]
  \node[font=\small\bfseries,coulExo] at (7,-1.3) {b) à e)};
  \quadrillage{-0.4}{-0.4}{14.4}{4.4}
  \pt{1,1}{above left}{A} \pt{6,1}{below left}{B} \pt{14,2}{above right}{C}
  \pt{4,4}{above}{D}
\end{tikzpicture}
\end{center}
\begin{correction}
\cas{a)}\par\nopagebreak
\explique{$\vec{u}$ : $4$ carreaux à droite et $2$ vers le haut.}{on part de l'origine :
d'abord horizontalement, puis verticalement}
\explique{$\vec{v}$ : $3$ carreaux à droite et $4$ vers le bas.}{on compte vers la pointe}
\explique{$\vec{w}$ : $3$ carreaux à gauche et $2$ vers le bas.}{la pointe est à gauche,
en bas}

\cas{b)}\par\nopagebreak
\begin{minipage}[t]{0.72\linewidth}
\explique{Depuis $A$, on compte $4$ carreaux à droite, puis $1$ vers le bas, et on place la
pointe.}{la flèche part de $A$}
\end{minipage}\hfill
\begin{minipage}[t]{0.26\linewidth}
\vspace{0pt}\centering
\begin{tikzpicture}[scale=0.45]
  \quadrillage{-0.4}{-0.4}{4.4}{1.4}
  \draw[vecteurC] (0,1) -- node[midway,above,font=\footnotesize]{$\vec{s}$} (4,0);
  \pt{0,1}{left}{A}
\end{tikzpicture}
\end{minipage}

\cas{c)}\par\nopagebreak
\begin{minipage}[t]{0.72\linewidth}
\explique{$\vec{u}$ fait $4$ carreaux à droite et $2$ vers le haut.}{lu en a)}
\explique{Depuis $B$, on fait ce déplacement et on place la pointe.}{la flèche part de $B$}
\end{minipage}\hfill
\begin{minipage}[t]{0.26\linewidth}
\vspace{0pt}\centering
\begin{tikzpicture}[scale=0.45]
  \quadrillage{-0.4}{-0.4}{4.4}{2.4}
  \draw[vecteur] (0,0) -- node[midway,above left,font=\footnotesize]{$\vec{u}$} (4,2);
  \pt{0,0}{left}{B}
\end{tikzpicture}
\end{minipage}

\cas{d)}\par\nopagebreak
\begin{minipage}[t]{0.72\linewidth}
\explique{La pointe doit être en $C$ : depuis $C$, on fait le déplacement à l'envers,
$4$ carreaux à gauche et $2$ vers le bas. On arrive à l'origine.}{$C$ est l'extrémité}
\explique{On trace la flèche de cette origine jusqu'à $C$.}{la pointe en $C$}
\end{minipage}\hfill
\begin{minipage}[t]{0.26\linewidth}
\vspace{0pt}\centering
\begin{tikzpicture}[scale=0.45]
  \quadrillage{-0.4}{-0.4}{4.4}{2.4}
  \draw[vecteur] (0,0) -- node[midway,above left,font=\footnotesize]{$\vec{u}$} (4,2);
  \pt{4,2}{right}{C}
\end{tikzpicture}
\end{minipage}

\cas{e)}\par\nopagebreak
\begin{minipage}[t]{0.72\linewidth}
\explique{$-\vec{u}$ a le déplacement contraire de celui de $\vec{u}$ : $4$ carreaux à
gauche et $2$ vers le bas.}{vecteurs opposés : même direction, même longueur, sens
contraire}
\explique{Depuis $D$, on fait ce déplacement et on place la pointe.}{la flèche part de $D$}
\end{minipage}\hfill
\begin{minipage}[t]{0.26\linewidth}
\vspace{0pt}\centering
\begin{tikzpicture}[scale=0.45]
  \quadrillage{-0.4}{-0.4}{4.4}{2.4}
  \draw[vecteurB] (4,2) -- node[midway,above left,font=\footnotesize]{$-\vec{u}$} (0,0);
  \pt{4,2}{right}{D}
\end{tikzpicture}
\end{minipage}
\end{correction}

\etape{À vous de jouer}
\begin{questions}[label=\alph*)]
  \item Lire le déplacement de chaque vecteur $\vec{p}$, $\vec{q}$ et $\vec{r}$.
  \item Tracer le vecteur $\vec{t}$ d'origine $D$ qui fait $3$ carreaux à droite et $2$
        vers le bas.
  \item Tracer un représentant de $\vec{p}$ d'origine $E$.
  \item Tracer un autre représentant de $\vec{p}$, d'extrémité $F$.
  \item Tracer un représentant de $-\vec{p}$ d'origine $G$.
\end{questions}
\begin{center}
\begin{tikzpicture}[scale=0.45]
  \node[font=\small\bfseries,coulExo] at (7,-1.3) {a)};
  \quadrillage{-0.4}{-0.4}{14.4}{4.4}
  \draw[vecteur] (0,0) -- node[midway,left,font=\footnotesize]{$\vec{p}$} (3,4);
  \draw[vecteurB] (5,3) -- node[midway,above right,font=\footnotesize]{$\vec{q}$} (9,1);
  \draw[vecteurC] (14,0) -- node[midway,below left,font=\footnotesize]{$\vec{r}$} (10,3);
\end{tikzpicture}
\quad
\begin{tikzpicture}[scale=0.45]
  \node[font=\small\bfseries,coulExo] at (7,-1.3) {b) à e)};
  \quadrillage{-0.4}{-0.4}{14.4}{4.4}
  \pt{1,3}{above left}{D} \pt{6,0}{below left}{E} \pt{14,4}{above right}{F}
  \pt{12,4}{above left}{G}
\end{tikzpicture}
\end{center}
\reponse{a) $\vec{p}$ : $3$ à droite et $4$ vers le haut ; $\vec{q}$ : $4$ à droite et $2$
vers le bas ; $\vec{r}$ : $4$ à gauche et $3$ vers le haut.\; d) l'origine est à $3$
carreaux à gauche et $4$ vers le bas de $F$.\; e)~$-\vec{p}$ : $3$ carreaux à gauche et
$4$ vers le bas, depuis $G$.}
\end{methodebox}

\afaire{exercices 5 à 10}

\partiecours{Translation}

\begin{definitionbox}[Définition 4 --- Translation de vecteur $\vec{u}$]
La \textbf{translation de vecteur $\vec{u}$}, notée $t_{\vec{u}}$, fait glisser chaque
point $M$ jusqu'au point $M'$ tel que $\vect{MM'}=\vec{u}$. Le point $M'$ est
l'\textbf{image} de $M$ par $t_{\vec{u}}$.
\end{definitionbox}

\begin{proprietebox}[Propriété 1]
Une translation conserve les longueurs, les angles et les aires : une figure et son
image sont superposables. Une droite a pour image une droite \textbf{parallèle}.
\end{proprietebox}

\begin{methodebox}[Méthode 3 --- Construire l'image d'un point, d'une figure]
On reporte le déplacement de $\vec{u}$ à partir du point à transformer ; pour une
figure, on transforme chaque sommet, puis on relie les images dans le même ordre.
\etapeexemples
Construire l'image $A'B'C'$ du triangle $ABC$ par la translation de vecteur $\vec{u}$.
\begin{center}
\begin{tikzpicture}[scale=0.6]
  \quadrillage{-0.4}{-0.4}{9.4}{5.4}
  \draw[vecteur] (4,5) -- (9,4);
  \node[font=\small,coulCourbe] at (6.5,4.85) {$\vec{u}$};
  \draw[fig] (0,2) -- (2,3) -- (1,4) -- cycle;
  \pt{0,2}{below left}{A} \pt{2,3}{right}{B} \pt{1,4}{above left}{C}

  \draw[fig,coulCourbeC] (5,1) -- (7,2) -- (6,3) -- cycle;
  \pt{5,1}{below left}{A'} \pt{7,2}{right}{B'} \pt{6,3}{above left}{C'}
  \draw[vecteur,line width=0.8pt] (0,2) -- node[midway,below,font=\footnotesize]{$\vect{AA'}$} (5,1);
  \draw[vecteur,line width=0.8pt] (2,3) -- node[midway,below,font=\footnotesize]{$\vect{BB'}$} (7,2);
  \draw[vecteur,line width=0.8pt] (1,4) -- node[midway,above,font=\footnotesize]{$\vect{CC'}$} (6,3);

\end{tikzpicture}
\end{center}
\begin{correction}
\cas{Résolution}\par\nopagebreak
\explique{$\vec{u}$ : $5$ carreaux à droite et $1$ vers le bas.}{on lit le déplacement}
\explique{On le reporte depuis $A$, depuis $B$ et depuis $C$ : $\vect{AA'}$,
$\vect{BB'}$ et $\vect{CC'}$ sont égaux à $\vec{u}$.}{chaque sommet glisse de la même façon}
\end{correction}

\etape{À vous de jouer}
Construire l'image $D'E'F'$ du triangle $DEF$ par la translation de vecteur $\vec{v}$.
\begin{center}
\begin{tikzpicture}[scale=0.6]
  \quadrillage{-0.4}{-0.4}{12.4}{5.4}
  \draw[vecteur] (8,0) -- node[midway,below right,font=\footnotesize]{$\vec{v}$} (12,2);
  \draw[fig] (0,1) -- (3,0) -- (2,3) -- cycle;
  \pt{0,1}{left}{D} \pt{3,0}{below}{E} \pt{2,3}{above left}{F}
\end{tikzpicture}
\end{center}
\reponse{$\vec{v}$ : $4$ carreaux à droite et $2$ vers le haut ; on le reporte depuis
$D$, $E$ et $F$.}
\end{methodebox}

\afaire{exercices 11 à 16}

\partiecours{Parallélogramme et égalité de vecteurs}

\begin{proprietebox}[Théorème 1 --- Parallélogramme]
Pour quatre points $A$, $B$, $C$, $D$ : \quad
$\vect{AB}=\vect{DC}$ \quad équivaut à \quad $ABCD$ est un parallélogramme
(éventuellement aplati, si les points sont alignés).
\begin{center}
\begin{tikzpicture}[scale=0.85]
  \draw[fig] (0,0) -- (3,0.6) -- (4,2.4) -- (1,1.8) -- cycle;
  \draw[vecteur] (0,0) -- node[midway,below,font=\footnotesize]{$\vect{AB}$} (3,0.6);
  \draw[vecteur] (1,1.8) -- node[midway,above,font=\footnotesize]{$\vect{DC}$} (4,2.4);
  \pt{0,0}{below left}{A} \pt{3,0.6}{below right}{B}
  \pt{4,2.4}{above right}{C} \pt{1,1.8}{above left}{D}
\end{tikzpicture}
\end{center}
\end{proprietebox}

\begin{remarquebox}[L'ordre des lettres]
On nomme un quadrilatère en faisant le tour de la figure : dans $ABCD$, on va de $A$ à
$B$, puis à $C$, puis à $D$. Les côtés opposés $[AB]$ et $[DC]$ donnent deux vecteurs de
même sens. Attention : $\vect{AB}=\vect{CD}$ donne le parallélogramme $ABDC$, et non
$ABCD$.
\end{remarquebox}

\begin{methodebox}[Méthode 4 --- Placer un point défini par une égalité de vecteurs]
On lit le déplacement du vecteur connu, puis on le reporte depuis l'origine imposée par
l'égalité. Le théorème 1 nomme alors le parallélogramme obtenu.
\etapeexemples
\begin{questions}
  \item Tracer le vecteur $\vect{EF}$, puis construire le point $H$ tel que
        $\vect{GH}=\vect{EF}$.
  \item Tracer le parallélogramme obtenu, puis le nommer.
  \item Démontrer que c'est un parallélogramme.
\end{questions}
\begin{center}
\begin{tikzpicture}[scale=0.6]
  \quadrillage{-0.4}{-0.4}{6.4}{4.4}
  \pt{0,1}{below left}{E} \pt{4,2}{below right}{F} \pt{1,3}{above left}{G}
\end{tikzpicture}
\end{center}
\begin{correction}
\cas{1\textdegree) et 2\textdegree)}\par\nopagebreak
\begin{center}
\begin{tikzpicture}[scale=0.6]
  \quadrillage{-0.4}{-0.4}{6.4}{4.4}
  \draw[fig] (0,1) -- (4,2) -- (5,4) -- (1,3) -- cycle;
  \draw[vecteur] (0,1) -- node[midway,below,font=\footnotesize]{$\vect{EF}$} (4,2);
  \draw[vecteurB] (1,3) -- node[midway,above,font=\footnotesize]{$\vect{GH}$} (5,4);
  \pt{0,1}{below left}{E} \pt{4,2}{below right}{F} \pt{1,3}{above left}{G}
  \pt{5,4}{above right}{H}
\end{tikzpicture}
\end{center}
\explique{$\vect{EF}$ : $4$ carreaux à droite et $1$ vers le haut.}{on lit le
déplacement}
\explique{Depuis $G$, on avance de $4$ à droite et $1$ vers le haut : c'est $H$.}{on le
reporte depuis $G$}
\explique{On obtient le parallélogramme $EFHG$ : on joint $E$, $F$, $H$ et $G$ dans cet
ordre.}{le tour de la figure}
\cas{3\textdegree)}\par\nopagebreak
\explique{On a construit $H$ tel que $\vect{EF}=\vect{GH}$.}{la construction}
\explique{D'après le théorème 1, $EFHG$ est un parallélogramme.}{$\vect{AB}=\vect{DC}$
donne $ABCD$}
\end{correction}

\etape{À vous de jouer}
\begin{questions}
  \item Tracer le vecteur $\vect{KL}$, puis construire le point $N$ tel que
        $\vect{MN}=\vect{KL}$.
  \item Tracer le parallélogramme obtenu, puis le nommer.
  \item Démontrer que c'est un parallélogramme.
\end{questions}
\begin{center}
\begin{tikzpicture}[scale=0.6]
  \quadrillage{-0.4}{-0.4}{6.4}{4.4}
  \pt{0,2}{left}{K} \pt{3,0}{below right}{L} \pt{2,4}{above left}{M}
\end{tikzpicture}
\end{center}
\reponse{1\textdegree) $N$ : depuis $M$, $3$ carreaux à droite et $2$ vers le bas.\;
2\textdegree) $KLNM$.\; 3\textdegree) $\vect{KL}=\vect{MN}$, donc $KLNM$ est un
parallélogramme (théorème 1).}
\end{methodebox}

\afaire{exercices 17 à 22}

\partiecours{Somme de deux vecteurs}

\begin{definitionbox}[Définition 5 --- Somme]
\begin{minipage}[t]{0.58\linewidth}\raggedright
Enchaîner la translation de vecteur $\vec{u}$ puis celle de vecteur $\vec{v}$ revient à
faire une seule translation, de vecteur $\vec{u}+\vec{v}$.

\smallskip
On trace les deux vecteurs \textbf{bout à bout} (l'origine du second sur l'extrémité du
premier) : la somme va du \textbf{départ du premier vecteur} à l'\textbf{arrivée du second
vecteur}.

\smallskip
L'ordre ne compte pas : $\vec{u}+\vec{v}=\vec{v}+\vec{u}$.

\smallskip
La \textbf{différence} est la somme avec l'opposé : $\vec{u}-\vec{v}=\vec{u}+(-\vec{v})$.
\end{minipage}\hfill
\begin{minipage}[t]{0.38\linewidth}
\vspace{0pt}\centering
\begin{tikzpicture}[scale=0.8]
  \draw[vecteur] (0,0) -- node[midway,below right,font=\footnotesize]{$\vec{u}$} (3,1);
  \draw[vecteurB] (3,1) -- node[midway,right,font=\footnotesize]{$\vec{v}$} (4,3);
  \draw[vecteurC] (0,0) -- node[midway,above left,font=\footnotesize]{$\vec{u}+\vec{v}$} (4,3);
  \node[below,font=\footnotesize] at (0,-0.1) {départ};
  \node[above,font=\footnotesize] at (4,3.1) {arrivée};
\end{tikzpicture}
\end{minipage}
\end{definitionbox}

\begin{methodebox}[Méthode 5 --- Construire une somme ou une différence]
On met les vecteurs \textbf{l'un à la suite de l'autre} : l'origine du second sur
l'extrémité du premier ; la somme va du départ du premier à l'arrivée du second.
\etapeexemples
\begin{questions}[label=\alph*)]
  \item Construire $\vec{w}=\vec{u}+\vec{v}$.
  \item Construire $\vec{z}=\vec{u}-\vec{v}$.
  \item Construire le représentant d'origine $A$ de $\vec{m}=\vect{AB}+\vect{CD}$.
  \item Construire le représentant d'origine $A$ de $\vec{n}=\vect{AB}+\vect{CA}$. À quel
        vecteur de la figure est-il égal ?
  \item Construire le point $D$ tel que $\vect{AD}=\vect{AB}+\vect{AC}$. Quelle est la
        nature de $ABDC$ ? Le tracer et le démontrer.
\end{questions}
\begin{center}
\begin{tikzpicture}[scale=0.5]
  \node[font=\small\bfseries,coulExo] at (3,-1.2) {a) et b)};
  \quadrillage{-0.4}{-0.4}{6.4}{3.4}
  \draw[vecteur] (0,1) -- node[midway,above left,font=\footnotesize]{$\vec{u}$} (3,2);
  \draw[vecteurB] (5,3) -- node[midway,right,font=\footnotesize]{$\vec{v}$} (6,1);
\end{tikzpicture}
\quad
\begin{tikzpicture}[scale=0.5]
  \node[font=\small\bfseries,coulExo] at (3.5,-1.2) {c) et d)};
  \quadrillage{-0.4}{-0.4}{7.4}{5.4}
  \pt{0,2}{left}{A} \pt{3,3}{below right}{B} \pt{1,5}{above}{C} \pt{5,4}{above}{D}
\end{tikzpicture}
\quad
\begin{tikzpicture}[scale=0.5]
  \node[font=\small\bfseries,coulExo] at (2.5,-1.2) {e)};
  \quadrillage{-0.4}{-0.4}{5.4}{4.4}
  \pt{0,0}{below left}{A} \pt{4,1}{below right}{B} \pt{1,3}{above left}{C}
\end{tikzpicture}
\end{center}
\begin{correction}
\cas{a)}\par\nopagebreak
\begin{minipage}[t]{0.72\linewidth}
\explique{On trace un représentant de $\vec{v}$ à la suite de $\vec{u}$.}{l'un à la suite
de l'autre}
\explique{$\vec{w}$ va du départ de $\vec{u}$ à l'arrivée de $\vec{v}$.}{départ, arrivée}
\end{minipage}\hfill
\begin{minipage}[t]{0.26\linewidth}
\vspace{0pt}\centering
\begin{tikzpicture}[scale=0.5]
  \quadrillage{-0.4}{-0.4}{4.4}{2.4}
  \draw[vecteur] (0,1) -- node[midway,above left,font=\footnotesize]{$\vec{u}$} (3,2);
  \draw[vecteurB] (3,2) -- node[midway,right,font=\footnotesize]{$\vec{v}$} (4,0);
  \draw[vecteurC] (0,1) -- node[midway,below left,font=\footnotesize]{$\vec{w}$} (4,0);
\end{tikzpicture}
\end{minipage}

\cas{b)}\par\nopagebreak
\begin{minipage}[t]{0.72\linewidth}
\explique{On trace un représentant de $-\vec{v}$ à la suite de
$\vec{u}$.}{$\vec{u}-\vec{v}=\vec{u}+(-\vec{v})$}
\explique{$\vec{z}$ va du départ de $\vec{u}$ à l'arrivée de $-\vec{v}$.}{départ, arrivée}
\end{minipage}\hfill
\begin{minipage}[t]{0.26\linewidth}
\vspace{0pt}\centering
\begin{tikzpicture}[scale=0.5]
  \quadrillage{-0.4}{-0.4}{3.4}{3.4}
  \draw[vecteur] (0,0) -- node[midway,below right,font=\footnotesize]{$\vec{u}$} (3,1);
  \draw[vecteurB] (3,1) -- node[midway,right,font=\footnotesize]{$-\vec{v}$} (2,3);
  \draw[vecteurC] (0,0) -- node[midway,above left,font=\footnotesize]{$\vec{z}$} (2,3);
\end{tikzpicture}
\end{minipage}

\cas{c)}\par\nopagebreak
\begin{minipage}[t]{0.72\linewidth}
\explique{On trace un représentant de $\vect{CD}$ à la suite de $\vect{AB}$, depuis
$B$.}{l'un à la suite de l'autre}
\explique{$\vec{m}$ va du départ de $\vect{AB}$ à l'arrivée de $\vect{CD}$.}{départ,
arrivée}
\end{minipage}\hfill
\begin{minipage}[t]{0.26\linewidth}
\vspace{0pt}\centering
\begin{tikzpicture}[scale=0.5]
  \quadrillage{-0.4}{-0.4}{7.4}{5.4}
  \draw[vecteur] (0,2) -- node[midway,above,font=\footnotesize]{$\vect{AB}$} (3,3);
  \draw[vecteurB] (1,5) -- node[midway,above,font=\footnotesize]{$\vect{CD}$} (5,4);
  \draw[vecteurB] (3,3) -- node[midway,above right,font=\footnotesize]{$\vect{CD}$} (7,2);
  \draw[vecteurC] (0,2) -- node[midway,below,font=\footnotesize]{$\vec{m}$} (7,2);
  \pt{0,2}{left}{A} \pt{3,3}{above}{B} \pt{1,5}{above}{C} \pt{5,4}{above}{D}
\end{tikzpicture}
\end{minipage}

\cas{d)}\par\nopagebreak
\begin{minipage}[t]{0.72\linewidth}
\explique{On trace un représentant de $\vect{CA}$ à la suite de $\vect{AB}$, depuis
$B$.}{l'un à la suite de l'autre}
\explique{$\vec{n}$ va du départ de $\vect{AB}$ à l'arrivée de $\vect{CA}$.}{départ,
arrivée}
\explique{$\vec{n}$ et $\vect{CB}$ (en pointillés) sont deux représentants du même
vecteur : $\vec{n}=\vect{CB}$.}{on compare}
\explique{On pouvait le prévoir : $\vec{n}=\vect{CA}+\vect{AB}$, deux vecteurs déjà bout à
bout, qui vont de $C$ à $B$.}{l'ordre ne compte pas}
\end{minipage}\hfill
\begin{minipage}[t]{0.26\linewidth}
\vspace{0pt}\centering
\begin{tikzpicture}[scale=0.5]
  \quadrillage{-0.4}{-0.4}{5.4}{5.4}
  \draw[aide] (1,5) -- (3,3);
  \draw[vecteur] (0,2) -- node[midway,above,font=\footnotesize]{$\vect{AB}$} (3,3);
  \draw[vecteurB] (1,5) -- node[midway,left,font=\footnotesize]{$\vect{CA}$} (0,2);
  \draw[vecteurB] (3,3) -- node[midway,right,font=\footnotesize]{$\vect{CA}$} (2,0);
  \draw[vecteurC] (0,2) -- node[midway,below left,font=\footnotesize]{$\vec{n}$} (2,0);
  \pt{0,2}{left}{A} \pt{3,3}{right}{B} \pt{1,5}{above}{C} \pt{5,4}{above}{D}
\end{tikzpicture}
\end{minipage}

\cas{e)}\par\nopagebreak
\begin{minipage}[t]{0.72\linewidth}
\explique{On trace un représentant de $\vect{AC}$ à la suite de $\vect{AB}$, depuis $B$ :
il arrive en $D$.}{l'un à la suite de l'autre}
\explique{On a construit $D$ tel que $\vect{AC}=\vect{BD}$.}{la construction}
\explique{D'après le théorème 1, $ACDB$, c'est-à-dire $ABDC$, est un
parallélogramme.}{$\vect{AB}=\vect{DC}$ donne $ABCD$}
\explique{Deux vecteurs de même origine : leur somme suit la diagonale du
parallélogramme.}{à retenir}
\end{minipage}\hfill
\begin{minipage}[t]{0.26\linewidth}
\vspace{0pt}\centering
\begin{tikzpicture}[scale=0.5]
  \quadrillage{-0.4}{-0.4}{5.4}{4.4}
  \draw[fig] (0,0) -- (4,1) -- (5,4) -- (1,3) -- cycle;
  \draw[vecteur] (0,0) -- node[midway,below right,font=\footnotesize]{$\vect{AB}$} (4,1);
  \draw[vecteurB] (0,0) -- node[midway,left,font=\footnotesize]{$\vect{AC}$} (1,3);
  \draw[vecteurB] (4,1) -- node[midway,right,font=\footnotesize]{$\vect{BD}$} (5,4);
  \draw[vecteurC] (0,0) -- node[midway,above left,font=\footnotesize]{$\vect{AD}$} (5,4);
  \pt{0,0}{below left}{A} \pt{4,1}{below right}{B} \pt{1,3}{above left}{C}
  \pt{5,4}{above right}{D}
\end{tikzpicture}
\end{minipage}
\end{correction}

\etape{À vous de jouer}
\begin{questions}[label=\alph*)]
  \item Construire $\vec{c}=\vec{a}+\vec{b}$.
  \item Construire $\vec{d}=\vec{a}-\vec{b}$.
  \item Construire le représentant d'origine $K$ de $\vec{p}=\vect{KL}+\vect{MN}$.
  \item Construire le représentant d'origine $K$ de $\vec{q}=\vect{KL}+\vect{MK}$. À quel
        vecteur de la figure est-il égal ?
  \item Construire le point $U$ tel que $\vect{RU}=\vect{RS}+\vect{RT}$. Quelle est la
        nature de $RSUT$ ? Le tracer et le démontrer.
\end{questions}
\begin{center}
\begin{tikzpicture}[scale=0.5]
  \node[font=\small\bfseries,coulExo] at (3,-1.2) {a) et b)};
  \quadrillage{-0.4}{-0.4}{6.4}{3.4}
  \draw[vecteur] (0,0) -- node[midway,left,font=\footnotesize]{$\vec{a}$} (1,3);
  \draw[vecteurB] (3,3) -- node[midway,above right,font=\footnotesize]{$\vec{b}$} (6,2);
\end{tikzpicture}
\quad
\begin{tikzpicture}[scale=0.5]
  \node[font=\small\bfseries,coulExo] at (3.5,-1.2) {c) et d)};
  \quadrillage{-0.4}{-0.4}{7.4}{5.4}
  \pt{0,2}{left}{K} \pt{4,3}{below right}{L} \pt{1,4}{above}{M} \pt{4,5}{above}{N}
\end{tikzpicture}
\quad
\begin{tikzpicture}[scale=0.5]
  \node[font=\small\bfseries,coulExo] at (2,-1.2) {e)};
  \quadrillage{-0.4}{-0.4}{4.4}{4.4}
  \pt{0,1}{left}{R} \pt{3,0}{below right}{S} \pt{1,4}{above left}{T}
\end{tikzpicture}
\end{center}
\reponse{a) $\vec{c}$ : $4$ à droite, $2$ vers le haut.\; b)~$\vec{d}$ : $2$ à gauche,
$4$ vers le haut.\; c)~$\vec{p}$ : $7$ à droite, $2$ vers le haut.\;
d)~$\vec{q}$ : $3$ à droite, $1$ vers le bas ; $\vec{q}=\vect{ML}$.\;
e)~$RSUT$ est un parallélogramme.}
\end{methodebox}

\afaire{exercices 23 à 27}

\partiecours{Multiplier un vecteur par un nombre réel}

\begin{definitionbox}[Définition 6 --- Produit par un nombre réel]
Pour un vecteur $\vec{u}$ non nul et un nombre réel $k$ non nul, le vecteur $k\vec{u}$ :
\begin{itemize}
  \item a la \textbf{même direction} que $\vec{u}$ ;
  \item a le \textbf{même sens} si $k>0$, le \textbf{sens contraire} si $k<0$ ;
  \item est $k$ fois plus long si $k>0$ (et $-k$ fois plus long si $k<0$).
\end{itemize}
Si $k=0$ ou $\vec{u}=\vec{0}$, alors $k\vec{u}=\vec{0}$. Par exemple,
$\vec{u}+\vec{u}+\vec{u}=3\vec{u}$ et $(-1)\vec{u}=-\vec{u}$.

\smallskip
Deux vecteurs de même direction, comme $\vec{u}$ et $k\vec{u}$, sont dits
\textbf{colinéaires}.
\end{definitionbox}

\begin{exemplebox}[Exemple]
\begin{center}
\begin{tikzpicture}[scale=0.4]
  \quadrillage{-0.4}{-0.4}{20.4}{3.4}
  \draw[vecteur] (0,1) -- node[midway,above,font=\footnotesize]{$\vec{u}$} (3,2);
  \draw[vecteurB] (4,0) -- node[midway,above left,font=\footnotesize]{$2\vec{u}$} (10,2);
  \draw[vecteurC] (20,3) -- node[midway,above left,font=\footnotesize]{$-3\vec{u}$} (11,0);
\end{tikzpicture}
\end{center}
$\vec{u}$ : $3$ carreaux à droite et $1$ vers le haut.

$2\vec{u}$ : même direction, même sens, deux fois plus long : $6$ carreaux à droite et $2$
vers le haut.

$-3\vec{u}$ : même direction, sens contraire, trois fois plus long : $9$ carreaux à gauche
et $3$ vers le bas.
\end{exemplebox}

\begin{methodebox}[Méthode 6 --- Construire $k\vec{u}$, puis une somme de produits]
On multiplie par $k$ les deux nombres de carreaux du déplacement de $\vec{u}$ ; si $k<0$,
on change le sens. Pour une somme de produits, on construit chaque produit, puis on les met
l'un à la suite de l'autre (méthode 5).
\etapeexemples
\begin{questions}[label=\alph*)]
  \item Construire $2\vec{u}$ et $-1{,}5\vec{u}$.
  \item Construire $2\vec{a}$ et $-3\vec{b}$, puis $\vec{c}=-3\vec{a}+2\vec{b}$.
  \item $ABC$ est un triangle. Placer le point $M$ tel que
        ${\vect{AM}=2\vect{AB}-\vect{AC}}$, puis le point $N$ tel que
        ${\vect{BN}=2\vect{BC}-\vect{AC}}$.
\end{questions}
\begin{center}
\begin{tikzpicture}[scale=0.45]
  \node[font=\small\bfseries,coulExo] at (2,-1.2) {a)};
  \quadrillage{-0.4}{-0.4}{4.4}{2.4}
  \draw[vecteur] (0,0) -- node[midway,above left,font=\footnotesize]{$\vec{u}$} (4,2);
\end{tikzpicture}
\quad
\begin{tikzpicture}[scale=0.45]
  \node[font=\small\bfseries,coulExo] at (2,-1.2) {b)};
  \quadrillage{-0.4}{-0.4}{4.4}{2.4}
  \draw[vecteur] (0,0) -- node[midway,below right,font=\footnotesize]{$\vec{a}$} (2,1);
  \draw[vecteurB] (3,0) -- node[midway,right,font=\footnotesize]{$\vec{b}$} (4,2);
\end{tikzpicture}
\quad
\begin{tikzpicture}[scale=0.45]
  \node[font=\small\bfseries,coulExo] at (2,-2.2) {c)};
  \quadrillage{-2.4}{-1.4}{6.4}{5.4}
  \draw[fig] (0,0) -- (3,1) -- (1,3) -- cycle;
  \pt{0,0}{below left}{A} \pt{3,1}{right}{B} \pt{1,3}{above}{C}
\end{tikzpicture}
\end{center}
\begin{correction}
\cas{a)}\par\nopagebreak
\begin{minipage}[t]{0.64\linewidth}
\explique{$\vec{u}$ : $4$ carreaux à droite et $2$ vers le haut.}{on lit le déplacement}
\explique{$2\vec{u}$ : $8$ carreaux à droite et $4$ vers le haut.}{même sens, deux fois
plus long}
\explique{$-1{,}5\vec{u}$ : $6$ carreaux à gauche et $3$ vers le bas.}{sens contraire,
une fois et demie plus long}
\end{minipage}\hfill
\begin{minipage}[t]{0.34\linewidth}
\vspace{0pt}\centering
\begin{tikzpicture}[scale=0.34]
  \quadrillage{-0.4}{-0.4}{15.4}{4.4}
  \draw[vecteurB] (0,0) -- node[midway,above left,font=\footnotesize]{$2\vec{u}$} (8,4);
  \draw[vecteurC] (15,4) -- node[midway,below right,font=\footnotesize]{$-1{,}5\vec{u}$} (9,1);
\end{tikzpicture}
\end{minipage}

\cas{b)}\par\nopagebreak
\begin{minipage}[t]{0.64\linewidth}
\explique{$2\vec{a}$ : $4$ carreaux à droite et $2$ vers le haut.}{$\vec{a}$ : $2$ à droite,
$1$ vers le haut}
\explique{$-3\vec{b}$ : $3$ carreaux à gauche et $6$ vers le bas.}{$\vec{b}$ : $1$ à
droite, $2$ vers le haut}
\end{minipage}\hfill
\begin{minipage}[t]{0.34\linewidth}
\vspace{0pt}\centering
\begin{tikzpicture}[scale=0.4]
  \quadrillage{-0.4}{-0.4}{9.4}{6.4}
  \draw[vecteur] (0,2) -- node[midway,above left,font=\footnotesize]{$2\vec{a}$} (4,4);
  \draw[vecteurB] (9,6) -- node[midway,right,font=\footnotesize]{$-3\vec{b}$} (6,0);
\end{tikzpicture}
\end{minipage}

\smallskip
\begin{minipage}[t]{0.64\linewidth}
\explique{On trace un représentant de $-3\vec{a}$ : $6$ carreaux à gauche et $3$ vers le
bas.}{sens contraire, trois fois plus long}
\explique{À la suite, un représentant de $2\vec{b}$ : $2$ carreaux à droite et $4$ vers le
haut.}{l'un à la suite de l'autre}
\explique{$\vec{c}$ va du départ de $-3\vec{a}$ à l'arrivée de $2\vec{b}$ : $4$ carreaux à
gauche et $1$ vers le haut.}{départ, arrivée}
\end{minipage}\hfill
\begin{minipage}[t]{0.34\linewidth}
\vspace{0pt}\centering
\begin{tikzpicture}[scale=0.4]
  \quadrillage{-0.4}{-0.4}{6.4}{4.4}
  \draw[vecteur] (6,3) -- node[midway,above,font=\footnotesize]{$-3\vec{a}$} (0,0);
  \draw[vecteurB] (0,0) -- node[midway,left,font=\footnotesize]{$2\vec{b}$} (2,4);
  \draw[vecteurC] (6,3) -- node[midway,above,font=\footnotesize]{$\vec{c}$} (2,4);
\end{tikzpicture}
\end{minipage}

\cas{c)}\par\nopagebreak
\begin{minipage}[t]{0.64\linewidth}
\explique{Depuis $A$, on trace un représentant de $2\vect{AB}$ : $6$ carreaux à droite et
$2$ vers le haut.}{$\vect{AB}$ : $3$ à droite, $1$ vers le haut}
\explique{À la suite, un représentant de $-\vect{AC}$ : $1$ carreau à gauche et $3$ vers
le bas.}{$\vect{AC}$ : $1$ à droite, $3$ vers le haut}
\explique{Il arrive en $M$ : $\vect{AM}=2\vect{AB}-\vect{AC}$.}{départ, arrivée}
\end{minipage}\hfill
\begin{minipage}[t]{0.34\linewidth}
\vspace{0pt}\centering
\begin{tikzpicture}[scale=0.4]
  \quadrillage{-0.4}{-1.4}{6.4}{3.4}
  \draw[fig] (0,0) -- (3,1) -- (1,3) -- cycle;
  \draw[vecteur] (0,0) -- node[pos=0.7,above,font=\footnotesize]{$2\vect{AB}$} (6,2);
  \draw[vecteurB] (6,2) -- node[midway,right,font=\footnotesize]{$-\vect{AC}$} (5,-1);
  \draw[vecteurC] (0,0) -- node[midway,below,font=\footnotesize]{$\vect{AM}$} (5,-1);
  \pt{0,0}{left}{A} \pt{3,1}{below}{B} \pt{1,3}{above}{C} \pt{5,-1}{below}{M}
\end{tikzpicture}
\end{minipage}

\smallskip
\begin{minipage}[t]{0.64\linewidth}
\explique{Depuis $B$, on trace un représentant de $2\vect{BC}$ : $4$ carreaux à gauche et
$4$ vers le haut.}{$\vect{BC}$ : $2$ à gauche, $2$ vers le haut}
\explique{À la suite, un représentant de $-\vect{AC}$ : $1$ carreau à gauche et $3$ vers
le bas.}{$\vect{AC}$ : $1$ à droite, $3$ vers le haut}
\explique{Il arrive en $N$ : $\vect{BN}=2\vect{BC}-\vect{AC}$.}{départ, arrivée}
\end{minipage}\hfill
\begin{minipage}[t]{0.34\linewidth}
\vspace{0pt}\centering
\begin{tikzpicture}[scale=0.4]
  \quadrillage{-2.4}{-0.4}{3.4}{5.4}
  \draw[fig] (0,0) -- (3,1) -- (1,3) -- cycle;
  \draw[vecteur] (3,1) -- node[pos=0.92,right,font=\footnotesize]{$2\vect{BC}$} (-1,5);
  \draw[vecteurB] (-1,5) -- node[midway,left,font=\footnotesize]{$-\vect{AC}$} (-2,2);
  \draw[vecteurC] (3,1) -- node[pos=0.3,below,font=\footnotesize]{$\vect{BN}$} (-2,2);
  \pt{0,0}{below}{A} \pt{3,1}{right}{B} \pt{1,3}{above}{C} \pt{-2,2}{left}{N}
\end{tikzpicture}
\end{minipage}
\end{correction}

\etape{À vous de jouer}
\begin{questions}[label=\alph*)]
  \item Construire $2\vec{v}$ et $-1{,}5\vec{v}$.
  \item Construire $2\vec{m}$ et $-3\vec{n}$, puis $\vec{p}=-3\vec{m}+2\vec{n}$.
  \item $RST$ est un triangle. Placer le point $P$ tel que
        ${\vect{RP}=2\vect{RS}-\vect{RT}}$, puis le point $U$ tel que
        ${\vect{SU}=2\vect{ST}-\vect{RT}}$.
\end{questions}
\begin{center}
\begin{tikzpicture}[scale=0.45]
  \node[font=\small\bfseries,coulExo] at (2,-1.2) {a)};
  \quadrillage{-0.4}{-0.4}{4.4}{2.4}
  \draw[vecteur] (0,2) -- node[midway,below left,font=\footnotesize]{$\vec{v}$} (4,0);
\end{tikzpicture}
\quad
\begin{tikzpicture}[scale=0.45]
  \node[font=\small\bfseries,coulExo] at (2.5,-1.2) {b)};
  \quadrillage{-0.4}{-0.4}{5.4}{2.4}
  \draw[vecteur] (0,2) -- node[midway,below left,font=\footnotesize]{$\vec{m}$} (2,1);
  \draw[vecteurB] (3,0) -- node[midway,above left,font=\footnotesize]{$\vec{n}$} (5,1);
\end{tikzpicture}
\quad
\begin{tikzpicture}[scale=0.45]
  \node[font=\small\bfseries,coulExo] at (2,-5.2) {c)};
  \quadrillage{-2.4}{-4.4}{6.4}{5.4}
  \draw[fig] (0,0) -- (3,-1) -- (1,2) -- cycle;
  \pt{0,0}{left}{R} \pt{3,-1}{right}{S} \pt{1,2}{above}{T}
\end{tikzpicture}
\end{center}
\reponse{a) $2\vec{v}$ : $8$ carreaux à droite et $4$ vers le bas ;\; $-1{,}5\vec{v}$ :
$6$ à gauche et $3$ vers le haut.\; b)~$2\vec{m}$ : $4$ à droite et $2$ vers le bas ;\;
$-3\vec{n}$ : $6$ à gauche et $3$ vers le bas ;\; $\vec{p}$ : $2$ à gauche et $5$ vers le
haut.\; c)~$\vect{RP}$ : $5$ à droite et $4$ vers le bas ;\; $\vect{SU}$ : $5$ à gauche et
$4$ vers le haut.}
\end{methodebox}

\afaire{exercices 28 à 31}

\partiecours{Relation de Chasles}

\begin{proprietebox}[Théorème 2 --- Relation de Chasles]
Pour tous points $A$, $B$ et $C$ : \quad $\vect{AB}+\vect{BC}=\vect{AC}$.

\begin{itemize}
  \item Aller de $A$ à $B$, puis de $B$ à $C$, c'est aller de $A$ à $C$.
  \item $\vect{AB}+\vect{BA}=\vect{AA}=\vec{0}$, donc $-\vect{AB}=\vect{BA}$.
  \item On peut passer par n'importe quel point de son choix :
        $\vect{AC}=\vect{AM}+\vect{MC}$.
\end{itemize}
\begin{center}
\begin{tikzpicture}[scale=0.85]
  \draw[vecteur] (0,0) -- node[midway,above left,font=\footnotesize]{$\vect{AB}$} (2.6,1.6);
  \draw[vecteurB] (2.6,1.6) -- node[midway,above right,font=\footnotesize]{$\vect{BC}$} (5.5,0.7);
  \draw[vecteurC] (0,0) -- node[midway,below,font=\footnotesize]{$\vect{AC}$} (5.5,0.7);
  \pt{0,0}{below left}{A} \pt{2.6,1.6}{above}{B} \pt{5.5,0.7}{below right}{C}
\end{tikzpicture}
\end{center}
\end{proprietebox}

\begin{methodebox}[Méthode 7 --- Simplifier une somme de vecteurs]
On remplace chaque $-\vect{XY}$ par $\vect{YX}$, puis on range les vecteurs pour que
l'extrémité de chacun soit l'origine du suivant.

Un même vecteur écrit deux fois se regroupe, comme $a+a=2a$ en calcul littéral :
aller de $A$ à $B$, puis encore de $A$ à $B$, c'est faire deux fois le trajet
(partie VI : $\vec{u}+\vec{u}=2\vec{u}$).
\[ \vect{AB}+\vect{AB}=2\vect{AB} \]
\etapeexemples
Simplifier :
\begin{questions}[label=\alph*)]
  \item $\vect{KL}+\vect{JK}$
  \item $\vect{AB}+\vect{DA}+\vect{CD}$
  \item $\vect{EF}-\vect{GF}+\vect{GE}$
  \item $\vect{BA}+\vect{AC}-\vect{CB}$
\end{questions}
\begin{correction}
\cas{a)}
\begin{calculs}
\vect{KL}+\vect{JK} &= \vect{JK}+\vect{KL} &\qquad& \rem{on range : $J\to K\to L$}\\
\vect{KL}+\vect{JK} &= \vect{JL} &\qquad& \rem{Chasles}
\end{calculs}
\cas{b)}
\begin{calculs}
\vect{AB}+\vect{DA}+\vect{CD} &= \vect{CD}+\vect{DA}+\vect{AB} &\qquad& \rem{on range : $C\to D\to A\to B$}\\
\vect{AB}+\vect{DA}+\vect{CD} &= \vect{CB} &\qquad& \rem{Chasles}
\end{calculs}
\cas{c)}
\begin{calculs}
\vect{EF}-\vect{GF}+\vect{GE} &= \vect{EF}+\vect{FG}+\vect{GE} &\qquad& \rem{$-\vect{GF}=\vect{FG}$}\\
\vect{EF}-\vect{GF}+\vect{GE} &= \vect{EE}=\vec{0} &\qquad& \rem{on revient au point de départ}
\end{calculs}
\cas{d)}
\begin{calculs}
\vect{BA}+\vect{AC}-\vect{CB} &= \vect{BA}+\vect{AC}+\vect{BC} &\qquad& \rem{$-\vect{CB}=\vect{BC}$}\\
\vect{BA}+\vect{AC}-\vect{CB} &= \vect{BC}+\vect{BC} &\qquad& \rem{Chasles : $\vect{BA}+\vect{AC}=\vect{BC}$}\\
\vect{BA}+\vect{AC}-\vect{CB} &= 2\vect{BC} &\qquad& \rem{on regroupe : $\vect{BC}$ deux fois}
\end{calculs}
\end{correction}

\begin{remarquebox}[Erreur à éviter]
Chasles ne s'applique que si la lettre d'arrivée du premier vecteur est la lettre de
départ du second : $\vect{AB}+\vect{CD}$ ne se simplifie pas.
\end{remarquebox}

\etape{À vous de jouer}
Simplifier :
\[
  \vect{UV}+\vect{VW} \qquad \vect{GH}+\vect{FG} \qquad \vect{MN}+\vect{NP}+\vect{PM}
  \qquad \vect{AC}+\vect{BA}+\vect{CB} \qquad \vect{EF}-\vect{GF}
\]
\[
  \vect{DE}+\vect{EF}-\vect{FD} \qquad \vect{KL}-\vect{ML}+\vect{KM}
\]
\reponse{$\vect{UW}$ ;\; $\vect{FH}$ ;\; $\vec{0}$ ;\; $\vect{BA}+\vect{AC}+\vect{CB}=\vec{0}$ ;\;
$\vect{EF}+\vect{FG}=\vect{EG}$ ;\; $\vect{DF}+\vect{DF}=2\vect{DF}$ ;\;
$\vect{KM}+\vect{KM}=2\vect{KM}$.}
\end{methodebox}

\afaire{exercices 32 à 37}

\partiecours{Démonstrations vectorielles}

\begin{proprietebox}[Propriété 2 --- Calculer avec des vecteurs]
On calcule avec les vecteurs comme avec les nombres en calcul littéral : pour tous
vecteurs $\vec{u}$, $\vec{v}$ et tous nombres réels $k$, $a$, $b$,
\begin{center}
$k\bigl(\vec{u}+\vec{v}\bigr)=k\vec{u}+k\vec{v}$ \qquad
$a\vec{u}+b\vec{u}=(a+b)\vec{u}$
\end{center}
comme $k(x+y)=kx+ky$ et $3x-x=2x$. Par exemple,
$2\bigl(\vect{AB}+\vect{BC}\bigr)=2\vect{AB}+2\vect{BC}$ et
$3\vect{AB}-\vect{AB}=2\vect{AB}$.
\end{proprietebox}

\begin{remarquebox}[Démontrer une égalité : le départ et le but]
Dans l'égalité à démontrer, le membre de \textbf{gauche} (avant le signe $=$) est le
\textbf{point de départ} ; le membre de \textbf{droite} (après le signe $=$) est le
\textbf{but}. On recopie le départ, on le transforme ligne par ligne, et on s'arrête
quand on a obtenu le but. Si le calcul semble plus facile dans l'autre sens, on peut
aussi partir du membre de droite et viser celui de gauche.

\smallskip
Tant que l'égalité n'est pas démontrée, on ne peut pas s'en servir : on ne fait
\textbf{jamais} passer un vecteur de l'autre côté du signe $=$, comme dans une équation.
\end{remarquebox}

\begin{methodebox}[Méthode 8 --- Démontrer une égalité de vecteurs]
On part d'un seul membre (celui de gauche, ou celui de droite si c'est plus facile),
recopié sur la première ligne. À chaque ligne, une seule
transformation : un opposé retourné ($-\vect{XY}=\vect{YX}$), la relation de Chasles, une
donnée de l'énoncé, un développement ou un regroupement (propriété 2). Pour passer par un
point avec Chasles, on choisit un point du but ou de l'énoncé. Chaque ligne commence par
le membre de départ, puis on conclut.
\etapeexemples
\begin{questions}[label=\alph*)]
  \item Démontrer que, pour tous points $A$, $B$, $C$ et $D$ :
        $\vect{AC}+\vect{BD}=\vect{AD}+\vect{BC}$.
  \item $ABC$ est un triangle et $D$ le point tel que $\vect{AD}=\vect{AB}+\vect{AC}$.
        Démontrer que $\vect{BD}=\vect{AC}$, puis en déduire la nature du quadrilatère
        $ACDB$.
  \item $ABC$ est un triangle. Exprimer en fonction de $\vect{AB}$ et $\vect{AC}$ les
        vecteurs ${\vec{u}=\vect{BC}}$ et ${\vec{v}=2\vect{CB}+\vect{BA}}$.
\end{questions}
\begin{correction}
\cas{a)}
\explique{Départ : $\vect{AC}+\vect{BD}$ ; but : $\vect{AD}+\vect{BC}$. Le but commence
par $\vect{AD}$ : on passe par $D$ dans $\vect{AC}$.}{on lit le but avant de calculer}
\begin{calculs}
\vect{AC}+\vect{BD} &= \vect{AD}+\vect{DC}+\vect{BD} &\qquad& \rem{Chasles : on passe par $D$}\\
\vect{AC}+\vect{BD} &= \vect{AD}+\vect{BD}+\vect{DC} &\qquad& \rem{on range : $B\to D\to C$}\\
\vect{AC}+\vect{BD} &= \res{\vect{AD}+\vect{BC}} &\qquad& \rem{Chasles : c'est le but}
\end{calculs}
\explique{Conclusion : pour tous points $A$, $B$, $C$, $D$,
$\vect{AC}+\vect{BD}=\vect{AD}+\vect{BC}$.}{on n'a jamais touché au membre de droite}
\cas{b)}
\begin{minipage}[t]{0.66\linewidth}\raggedright
\explique{Départ : $\vect{BD}$ ; but : $\vect{AC}$. L'énoncé parle de $\vect{AD}$ : on
passe par $A$.}{on passe par un point de l'énoncé}
\begin{calculs}
\vect{BD} &= \vect{BA}+\vect{AD} &\qquad& \rem{Chasles : on passe par $A$}\\
\vect{BD} &= \vect{BA}+\vect{AB}+\vect{AC} &\qquad& \rem{on remplace $\vect{AD}$ (énoncé)}\\
\vect{BD} &= \res{\vect{AC}} &\qquad& \rem{$\vect{BA}+\vect{AB}=\vect{BB}=\vec{0}$}
\end{calculs}
\explique{Conclusion : $\vect{AC}=\vect{BD}$, donc \res{$ACDB$ est un
parallélogramme}.}{théorème 1}
\end{minipage}\hfill
\begin{minipage}[t]{0.3\linewidth}\vspace{0pt}\centering
\begin{tikzpicture}[scale=0.7]
  \draw[fig] (0,0) -- (1,2) -- (4,2.5) -- (3,0.5) -- cycle;
  \draw[vecteur] (0,0) -- node[midway,left,font=\footnotesize]{$\vect{AC}$} (1,2);
  \draw[vecteur] (3,0.5) -- node[midway,right,font=\footnotesize]{$\vect{BD}$} (4,2.5);
  \pt{0,0}{below left}{A} \pt{3,0.5}{below right}{B}
  \pt{1,2}{above left}{C} \pt{4,2.5}{above right}{D}
\end{tikzpicture}
\end{minipage}
\cas{c)}
\explique{Les deux vecteurs autorisés partent de $A$ : on passe par $A$.}{le point
commun}
\begin{calculs}
\vec{u} &= \vect{BC} &\qquad& \rem{on recopie}\\
\vec{u} &= \vect{BA}+\vect{AC} &\qquad& \rem{Chasles : on passe par $A$}\\
\vec{u} &= \res{-\vect{AB}+\vect{AC}} &\qquad& \rem{$\vect{BA}=-\vect{AB}$}
\end{calculs}
\begin{calculs}
\vec{v} &= 2\vect{CB}+\vect{BA} &\qquad& \rem{on recopie}\\
\vec{v} &= 2\bigl(\vect{CA}+\vect{AB}\bigr)+\vect{BA} &\qquad& \rem{Chasles : on passe par $A$}\\
\vec{v} &= 2\vect{CA}+2\vect{AB}+\vect{BA} &\qquad& \rem{on développe (propriété 2)}\\
\vec{v} &= -2\vect{AC}+2\vect{AB}-\vect{AB} &\qquad& \rem{on retourne $\vect{CA}$ et $\vect{BA}$}\\
\vec{v} &= \res{\vect{AB}-2\vect{AC}} &\qquad& \rem{on regroupe : $2-1=1$}
\end{calculs}
\end{correction}

\begin{remarquebox}[Erreur à éviter]
Écrire l'égalité à démontrer en première ligne, puis « passer » des vecteurs de l'autre
côté : on se sert de ce qu'il faut démontrer. On part d'un seul membre.
\end{remarquebox}

\etape{À vous de jouer}
\begin{questions}[label=\alph*)]
  \item Démontrer que, pour tous points $E$, $F$, $G$ et $H$ :
        $\vect{EG}+\vect{FH}=\vect{EH}+\vect{FG}$.
  \item $EFG$ est un triangle et $H$ le point tel que $\vect{EH}=\vect{EF}+\vect{EG}$.
        Démontrer que $\vect{FH}=\vect{EG}$, puis en déduire la nature du quadrilatère
        $EGHF$.
  \item $EFG$ est un triangle. Exprimer en fonction de $\vect{EF}$ et $\vect{EG}$ les
        vecteurs ${\vec{u}=\vect{GF}}$ et ${\vec{v}=3\vect{FG}+\vect{GE}}$.
\end{questions}
\reponse{a) $\vect{EG}+\vect{FH}=\vect{EH}+\vect{HG}+\vect{FH}=\vect{EH}+\vect{FG}$.\;
b)~$\vect{FH}=\vect{FE}+\vect{EF}+\vect{EG}=\vect{EG}$ : $EGHF$ est un parallélogramme.\;
c)~$\vec{u}=\vect{EF}-\vect{EG}$ ;\; $\vec{v}=-3\vect{EF}+2\vect{EG}$.}
\end{methodebox}

\afaire{exercices 38 à 40}

\partiecours{Milieu d'un segment}

\begin{proprietebox}[Théorème 3 --- Milieu]
$I$ est le milieu de $[AB]$ \quad équivaut à \quad $\vect{AI}=\vect{IB}$.
\begin{center}
\begin{tikzpicture}[scale=0.85]
  \draw[vecteur] (0,0) -- node[midway,above left,font=\footnotesize]{$\vect{AI}$} (2.5,0.8);
  \draw[vecteur] (2.5,0.8) -- node[midway,above left,font=\footnotesize]{$\vect{IB}$} (5,1.6);
  \pt{0,0}{below left}{A} \pt{2.5,0.8}{below right}{I} \pt{5,1.6}{below right}{B}
\end{tikzpicture}
\end{center}
Pour démontrer que $M$ est le milieu de $[XY]$, il suffit donc d'obtenir
$\vect{XM}=\vect{MY}$.

\smallskip
On l'écrit aussi $\vect{IA}=-\vect{IB}$, ou $\vect{IA}+\vect{IB}=\vec{0}$. Rappel du
collège : un quadrilatère est un parallélogramme si, et seulement si, ses diagonales ont
le même milieu.
\end{proprietebox}

\begin{methodebox}[Méthode 9 --- Démontrer qu'un point est un milieu]
On vise l'égalité $\vect{XM}=\vect{MY}$. Un parallélogramme de l'énoncé se traduit par
une égalité de vecteurs (théorème 1) ; une somme se transforme en passant par un point
bien choisi (relation de Chasles), en partant d'un seul membre (méthode 8).
\etapeexemples
\begin{questions}[label=\alph*)]
  \item $ABCD$ et $BCFD$ sont deux parallélogrammes. Démontrer que $D$ est le milieu de
        $[AF]$.
  \item $ABC$ est un triangle et $D$ le point tel que $\vect{AD}=\vect{AB}+\vect{CB}$.
        Démontrer que $B$ est le milieu de $[CD]$.
\end{questions}
\begin{center}
\begin{tikzpicture}[scale=0.75]
  \node[font=\small\bfseries,coulExo] at (2,-0.8) {a)};
  \coordinate (A) at (0,0);
  \coordinate (B) at (3,0);
  \coordinate (C) at (4,1.6);
  \coordinate (D) at (1,1.6);
  \coordinate (F) at (2,3.2);
  \draw[fig] (A) -- (B) -- (C) -- (D) -- cycle;
  \draw[fig,black!45] (B) -- (C) -- (F) -- (D) -- cycle;
  \pt{A}{below left}{A} \pt{B}{below right}{B} \pt{C}{right}{C}
  \pt{D}{left}{D} \pt{F}{above}{F}
\end{tikzpicture}
\qquad
\begin{tikzpicture}[scale=0.75]
  \node[font=\small\bfseries,coulExo] at (2.5,-1.8) {b)};
  \draw[fig] (0,0) -- (3,1) -- (1,3) -- cycle;
  \draw[aide] (1,3) -- (5,-1);
  \pt{0,0}{below left}{A} \pt{3,1}{above right}{B} \pt{1,3}{above}{C}
  \pt{5,-1}{below right}{D}
\end{tikzpicture}
\end{center}
\begin{correction}
\cas{a)}\par\nopagebreak
\explique{$ABCD$ est un parallélogramme, donc $\vect{AD}=\vect{BC}$.}{côtés opposés}
\explique{$BCFD$ est un parallélogramme, donc $\vect{BC}=\vect{DF}$.}{côtés opposés}
\explique{Ainsi $\vect{AD}=\vect{DF}$.}{on enchaîne les deux égalités}
\explique{Conclusion : $D$ est le milieu de $[AF]$.}{théorème 3}
\cas{b)}
\explique{On part de $\vect{BD}$ et on passe par $A$, qui figure dans l'énoncé.}{on vise
$\vect{CB}=\vect{BD}$}
\begin{calculs}
\vect{BD} &= \vect{BA}+\vect{AD} &\qquad& \rem{Chasles : on passe par $A$}\\
\vect{BD} &= \vect{BA}+\vect{AB}+\vect{CB} &\qquad& \rem{on remplace $\vect{AD}$ (énoncé)}\\
\vect{BD} &= \vect{CB} &\qquad& \rem{$\vect{BA}+\vect{AB}=\vect{BB}=\vec{0}$}
\end{calculs}
\explique{Conclusion : $\vect{CB}=\vect{BD}$, donc $B$ est le milieu de $[CD]$.}{théorème 3}
\end{correction}

\etape{À vous de jouer}
\begin{questions}[label=\alph*)]
  \item $ABCD$ et $ABEC$ sont deux parallélogrammes. Démontrer que $C$ est le milieu de
        $[DE]$.
  \item $EFG$ est un triangle et $H$ le point tel que $\vect{EH}=\vect{EF}+\vect{GF}$.
        Démontrer que $F$ est le milieu de $[GH]$.
\end{questions}
\begin{center}
\begin{tikzpicture}[scale=0.75]
  \node[font=\small\bfseries,coulExo] at (3.5,-0.8) {a)};
  \coordinate (A) at (0,0);
  \coordinate (B) at (3,0);
  \coordinate (C) at (4,1.6);
  \coordinate (D) at (1,1.6);
  \coordinate (E) at (7,1.6);
  \draw[fig] (A) -- (B) -- (C) -- (D) -- cycle;
  \draw[fig,black!45] (A) -- (B) -- (E) -- (C) -- cycle;
  \pt{A}{below left}{A} \pt{B}{below right}{B} \pt{C}{above}{C}
  \pt{D}{above left}{D} \pt{E}{above right}{E}
\end{tikzpicture}
\qquad
\begin{tikzpicture}[scale=0.75]
  \node[font=\small\bfseries,coulExo] at (3,-1.8) {b)};
  \draw[fig] (0,0) -- (4,1) -- (2,3) -- cycle;
  \draw[aide] (2,3) -- (6,-1);
  \pt{0,0}{below left}{E} \pt{4,1}{above right}{F} \pt{2,3}{above}{G}
  \pt{6,-1}{below right}{H}
\end{tikzpicture}
\end{center}
\reponse{a) $\vect{DC}=\vect{AB}$ et $\vect{AB}=\vect{CE}$, donc $\vect{DC}=\vect{CE}$ : $C$
est le milieu de $[DE]$.\; b)~$\vect{FH}=\vect{FE}+\vect{EH}=\vect{GF}$ : $F$ est le milieu
de $[GH]$.}
\end{methodebox}

\afaire{exercices 41 à 45}

\end{document}
